\scanpage{614}\mspage{2}
% archive sheet number: 613

sont satisfaites simultanément. Quand aucune confusion n'est à
craindre, il m'arrivera de dire \inserted{simplement} que $\Phi$
est Q-clos, ou que $\Psi$ est Q-clos, sans spécifier si c'est à
g. ou à dr. Mais pour un couple Q-orthogonal $(\Phi,\Psi)$ (ou je
préfère l'écrire $\Phi \mathrel{\tikz[baseline=-.45ex,x=1em,y=1ex]{\draw (0,0)--(2.15,0) (0.32,-0.5)--(0,0) (1.83,0.5)--(2.15,0);}} \Psi$), Q-clos signifiera
tjrs Q-clos à gauche et à droite.

Quand un couple Q-orthogonal $\Phi \mathrel{\tikz[baseline=-.45ex,x=1em,y=1ex]{\draw (0,0)--(2.15,0) (0.32,-0.5)--(0,0) (1.83,0.5)--(2.15,0);}} \Psi$ a la
propriété de factorisation de Quillen, je l'indique par la
\deleted{\unsure{ligne}} barre verticale croisant l'horizontale, ainsi
\[
 \Phi \mathrel{\tikz[baseline=-.45ex,x=1em,y=1ex]{\draw (0,0)--(2.6,0) (0.32,-0.5)--(0,0) (2.28,0.5)--(2.6,0) (1.3,-.52)--(1.3,.52);}} \Psi ,\tag*{(1.4)}
\]
et on définit de $=$ le sens des graphismes
\[
\left\{\begin{aligned}
 &\Phi \mathrel{\tikz[baseline=-.45ex,x=1em,y=1ex]{\draw (0,0)--(3.5,0) (0.32,-0.5)--(0,0) (3.18,0.5)--(3.5,0) (1.75,-.52)--(1.75,.52);}} \Psi\\
 &\Phi \unsure{\mathrel{\tikz[baseline=-.45ex,x=1em,y=1ex]{\draw (0,0)--(3.5,0) (0.32,0.5)--(0,0) (0.32,-0.5)--(0,0) (3.18,0.5)--(3.5,0) (3.31,-.25)--(3.5,0) (1.75,-.52)--(1.75,.52);}}|\mathrel{\tikz[baseline=-.45ex,x=1em,y=1ex]{\draw (0,0)--(3.5,0) (0.32,0.5)--(0,0) (0.32,-0.5)--(0,0) (3.18,0.5)--(3.5,0) (1.75,-.52)--(1.75,.52);}}} \Psi\\
 &\Phi \mathrel{\tikz[baseline=-.45ex,x=1em,y=1ex]{\draw (0,0)--(3.5,0) (0.32,-0.5)--(0,0) (3.18,0.5)--(3.5,0) (3.18,-0.5)--(3.5,0) (1.75,-.52)--(1.75,.52);}} \Psi\\
 &\Phi \mathrel{\tikz[baseline=-.45ex,x=1em,y=1ex]{\draw (0,0)--(3.5,0) (0.32,0.5)--(0,0) (0.32,-0.5)--(0,0) (3.18,0.5)--(3.5,0) (3.18,-0.5)--(3.5,0) (1.75,-.52)--(1.75,.52);}} \Psi .
\end{aligned}\right.\tag*{(1.5)}
\]

Un couple Q-\uline{orthogonal} Q-\uline{factorisant}
\deleted{$\Phi,\Psi$} $\Phi \mathrel{\tikz[baseline=-.45ex,x=1em,y=1ex]{\draw (0,0)--(2.6,0) (0.32,-0.5)--(0,0) (2.28,0.5)--(2.6,0) (1.3,-.52)--(1.3,.52);}} \Psi$ est appelé un
\uline{précouple de Quillen}. S'il est de plus Q-clos, i.e.
\deleted{\unsure{lié}} si on a $\Phi \mathrel{\tikz[baseline=-.45ex,x=1em,y=1ex]{\draw (0,0)--(2.6,0) (0.32,0.5)--(0,0) (0.32,-0.5)--(0,0) (2.28,0.5)--(2.6,0) (2.28,-0.5)--(2.6,0) (1.3,-.52)--(1.3,.52);}} \Psi$, on dit que
c'est un \uline{couple de Quillen}.

$\Phi$ est dit \uline{Q-saturé} \uline{à gauche} s'il
\inserted{contient les iso et} est stable par cochangement de base,
par \inserted{(petites)} compositions transfinies, et par facteurs
directs. On \unsure{définit} dualement \deleted{la qualité
\unsure{de} $\Psi$} \deleted{$\unsure{\Phi}$}$\unsure{\Psi|\Phi}$ est dit
\uline{Q-saturé à droite},
